% n9_1_v0.1_near_diagonal_decomp.tex
%
% N9.1 — Decomposition of Π_{j,x₀,k} − Π_j into near-diagonal
%         commutator and paraproduct pieces for j ∈ [k, k+J]
%
% NS lane analysis note — DO NOT reference Collatz deliverables.
% Part of the D9 envelope framework insertion programme.
%
% Author: DM3-lab · 2026
% -----------------------------------------------------------------------

\documentclass[11pt,a4paper]{article}

\usepackage{amsmath,amssymb,amsthm}
\usepackage{mathtools}
\usepackage{hyperref}
\usepackage{geometry}
\geometry{margin=2.5cm}

% -----------------------------------------------------------------------
% Theorem environments
% -----------------------------------------------------------------------
\newtheorem{theorem}{Theorem}[section]
\newtheorem{proposition}[theorem]{Proposition}
\newtheorem{lemma}[theorem]{Lemma}
\newtheorem{corollary}[theorem]{Corollary}
\newtheorem{definition}[theorem]{Definition}
\newtheorem{remark}[theorem]{Remark}

% -----------------------------------------------------------------------
% Notation macros
% -----------------------------------------------------------------------
\newcommand{\R}{\mathbb{R}}
\newcommand{\Z}{\mathbb{Z}}
\newcommand{\N}{\mathbb{N}}
\newcommand{\Pij}{\Pi_j}
\newcommand{\Pijxk}{\Pi_{j,x_0,k}}
\newcommand{\diff}{\Pijxk - \Pij}
\newcommand{\psi}{\psi}
\newcommand{\supp}{\operatorname{supp}}
\newcommand{\norm}[2]{\left\|#1\right\|_{#2}}
\newcommand{\abs}[1]{\left|#1\right|}
\newcommand{\ip}[2]{\langle #1,\, #2 \rangle}
\newcommand{\Lp}[1]{L^{#1}(\R^n)}
\newcommand{\Sob}[2]{H^{#1}(\R^n)}
\newcommand{\para}{\mathsf{P}}     % paraproduct
\newcommand{\comm}{\mathsf{C}}     % commutator
\newcommand{\res}{\mathsf{R}}      % resonant / remainder
\newcommand{\BHT}{\mathrm{BHT}}

% -----------------------------------------------------------------------
\title{%
  \textbf{N9.1 — Decomposition of $\Pi_{j,x_0,k} - \Pi_j$}\\[4pt]
  \large Explicit Near-Diagonal Commutator/Paraproduct Pieces\\
  for $j \in [k,\, k+J]$
}
\author{DM3-lab NS lane}
\date{Version 0.1 — 2026}
% -----------------------------------------------------------------------

\begin{document}

\maketitle
\tableofcontents

% -----------------------------------------------------------------------
\section{Overview and Goal}
% -----------------------------------------------------------------------

This note provides an explicit, finite decomposition of the near-diagonal
operator difference
\[
  \delta\Pi_{j,x_0,k} \;:=\; \Pi_{j,x_0,k} - \Pi_j
\]
into commutator and paraproduct pieces, valid for the scale range
$j \in [k,\, k+J]$ with $J \geq 1$ fixed.  All constants are tracked
explicitly with their dependence on $J$, $n$ (spatial dimension), and the
Schwartz seminorms of the underlying cutoff $\psi$.

\medskip\noindent
\textbf{Standing context.}
This decomposition is a building block of the \emph{D9 envelope framework}
(see Section~\ref{sec:d9link}).  The operator $\Pi_j$ is the standard
Littlewood--Paley projector at dyadic frequency $2^j$, and $\Pi_{j,x_0,k}$
is a spatially localised variant centered at $x_0 \in \R^n$ with spatial
cutoff at scale $2^{-k}$.

% -----------------------------------------------------------------------
\section{Setup and Definitions}
\label{sec:setup}
% -----------------------------------------------------------------------

\subsection{Littlewood--Paley projectors}

Fix a smooth radial function $\varphi \in C^\infty_c(\R^n)$ with
\[
  \supp\varphi \;\subset\; \{\xi : |\xi| \leq 2\}, \qquad
  \varphi(\xi) = 1 \text{ for } |\xi| \leq 1,
\]
and let $\psi(\xi) := \varphi(\xi) - \varphi(2\xi)$ be the associated
annular cut, so $\supp\psi \subset \{1/2 \leq |\xi| \leq 2\}$.
Define
\[
  \psi_j(\xi) := \psi(2^{-j}\xi), \qquad
  \widehat{(\Pi_j f)}(\xi) := \psi_j(\xi)\,\hat{f}(\xi).
\]
Then $\Pi_j$ is an $L^p$-bounded Fourier multiplier for all $1 < p < \infty$,
with
\[
  \norm{\Pi_j f}{L^p} \;\leq\; C_p\,\norm{f}{L^p}, \qquad
  C_p \text{ independent of } j.
\]

\subsection{Localised projector $\Pi_{j,x_0,k}$}

Let $\chi \in C^\infty_c(\R^n)$ satisfy $\chi(x) = 1$ for $|x| \leq 1$
and $\supp\chi \subset \{|x| \leq 2\}$.  Set
\[
  \chi_{x_0,k}(x) \;:=\; \chi\!\left(2^k(x - x_0)\right),
\]
a bump function of unit scale $2^{-k}$ centred at $x_0$.
Define the \emph{localised Littlewood--Paley projector}
\[
  \Pi_{j,x_0,k} f \;:=\; \chi_{x_0,k}\,\Pi_j(\chi_{x_0,k}\,f).
\]
This operator is the natural object appearing in the D9 envelope analysis,
where one needs control of frequency-localised energy restricted to a spatial
ball of radius $2^{-k}$.

\subsection{The near-diagonal range}

We say $j \in [k, k+J]$ is the \emph{near-diagonal} or
\emph{near-resonant} range: the spatial localisation scale $2^{-k}$ is
comparable (up to a factor $2^J$) to the frequency scale $2^{-j}$.
For $j \gg k+J$ the projector $\Pi_j$ is effectively supported well below
the localisation scale, and for $j \ll k$ the frequency and spatial scales
are far apart; both regimes admit easier estimates (see Remark~\ref{rem:offdiag}).

% -----------------------------------------------------------------------
\section{The Decomposition}
\label{sec:decomp}
% -----------------------------------------------------------------------

\subsection{Expansion of the difference}

Write $m := \chi_{x_0,k}$ for brevity.  Then
\[
  \Pi_{j,x_0,k} f - \Pi_j f
  \;=\; m\,\Pi_j(mf) - \Pi_j f
  \;=\; m\,\Pi_j(mf) - m\,\Pi_j f + m\,\Pi_j f - \Pi_j f.
\]
Rearranging:
\begin{equation}
  \label{eq:firstsplit}
  \delta\Pi_{j,x_0,k} f
  \;=\; m\bigl[\Pi_j(mf) - \Pi_j f\bigr]
        + (m-1)\,\Pi_j f.
\end{equation}
We decompose each bracket using the Bony paraproduct splitting.

\subsection{Bony paraproduct splitting}

For two functions $a, f$ write the bilinear product $af$ as
\[
  af \;=\; \para_a f + \para_f a + \res(a,f),
\]
where
\begin{align*}
  \para_a f &:= \sum_{\ell} S_{\ell-3} a \cdot \Pi_\ell f
             & &\text{(paraproduct: low freq.\ of $a$ × high freq.\ of $f$)},\\
  \para_f a &:= \sum_{\ell} S_{\ell-3} f \cdot \Pi_\ell a
             & &\text{(paraproduct: low freq.\ of $f$ × high freq.\ of $a$)},\\
  \res(a,f) &:= \sum_{\ell} \Pi_\ell a \cdot \widetilde\Pi_\ell f
             & &\text{(resonant/diagonal: comparable frequencies)},
\end{align*}
with $S_j := \sum_{\ell \leq j} \Pi_\ell$ and $\widetilde\Pi_j$ a slightly
fattened version of $\Pi_j$ covering $|\xi| \sim 2^j$.

\subsection{Decomposition into four pieces}

Apply the Bony splitting to $(m-1)f$, i.e.,
$(m-1)f = \para_{(m-1)} f + \para_f(m-1) + \res(m-1,f)$,
and substitute into \eqref{eq:firstsplit}.  Denoting the commutator
\[
  \comm(m,\Pi_j)f \;:=\; m\,\Pi_j(m f) - m\,\Pi_j f
  \;=\; m\,\Pi_j\bigl(\para_{(m-1)} f
        + \para_f(m-1) + \res(m-1,f)\bigr),
\]
the final \emph{four-term} decomposition is:
\begin{equation}
  \label{eq:main}
  \boxed{
  \delta\Pi_{j,x_0,k} f
  \;=\;
  \underbrace{m\,\Pi_j\!\bigl(\para_{(m-1)} f\bigr)}_{\displaystyle T_1}
  +\underbrace{m\,\Pi_j\!\bigl(\para_f(m-1)\bigr)}_{\displaystyle T_2}
  +\underbrace{m\,\Pi_j\!\bigl(\res(m-1,f)\bigr)}_{\displaystyle T_3}
  +\underbrace{(m-1)\,\Pi_j f}_{\displaystyle T_4}
  }
\end{equation}
where $T_1$--$T_3$ constitute the commutator $\comm(m,\Pi_j)f$ and $T_4$
is the direct localisation error.  The four interaction types are:

\begin{description}
  \item[Low--high ($T_1$):]
    The low-frequency part of $(m-1)$ (via $S_{\ell-3}(m-1)$) multiplies
    the high-frequency part of $f$ ($\Pi_\ell f$), projected back to scale
    $2^j$.  This is the standard paraproduct piece.

  \item[High--low ($T_2$):]
    The high-frequency part of $(m-1)$ ($\Pi_\ell(m-1)$ for $\ell \sim j$)
    multiplies the low-frequency part of $f$ ($S_{\ell-3}f$).  Since
    $m = \chi_{x_0,k}$ is smooth at scale $2^{-k}$, its spectrum is
    concentrated at $|\xi| \lesssim 2^k$, making this term sensitive to
    the near-diagonal range $j \in [k, k+J]$.

  \item[Resonant/diagonal ($T_3$):]
    Both $(m-1)$ and $f$ contribute at comparable frequency scales $\ell \sim j$.
    This is the hardest piece and is the primary source of the $J$-dependence.

  \item[Direct localisation error ($T_4$):]
    The cutoff mismatch $(m-1)\Pi_j f$ is supported away from the ball
    $B(x_0, 2^{-k})$ and is controlled by the $L^2$ norm of $f$ outside
    the support region.
\end{description}

% -----------------------------------------------------------------------
\section{Estimates for Each Piece}
\label{sec:estimates}
% -----------------------------------------------------------------------

Throughout this section, all constants $C$ may depend on $n$, $p$,
$\norm{\psi}{C^N}$, and $\norm{\chi}{C^N}$ for a sufficiently large $N$,
but are \emph{uniform in $j$, $k$, $x_0$} unless stated otherwise.

\subsection{Piece $T_4$: low-frequency contribution of $(m-1)$}

\begin{proposition}[Estimate for $T_4$]
\label{prop:T4}
For all $j \in [k, k+J]$ and $f \in L^2(\R^n)$,
\[
  \norm{T_4 f}{L^2}
  \;=\; \norm{(m-1)\Pi_j f}{L^2}
  \;\leq\; \norm{\Pi_j f}{L^2(\R^n \setminus B(x_0,2^{-k}))}
  \;\leq\; \norm{f}{L^2(\R^n)}.
\]
The first inequality holds because $m = 1$ on $B(x_0, 2^{-k})$ so
$(m-1)\Pi_j f$ is supported in $\{|x-x_0| \geq 2^{-k}\}$, giving the
pointwise bound
\[
  \abs{(m(x)-1)\Pi_j f(x)}
  \;\leq\;
  \mathbf{1}_{|x-x_0| \geq 2^{-k}}\,|\Pi_j f(x)|.
\]
The second inequality follows from $\norm{\Pi_j}{L^2 \to L^2} \leq C$
(with $C$ independent of $j$).  Hence
\[
  \norm{T_4 f}{L^2} \;\leq\; C\,\norm{f}{L^2(\R^n)},
\]
with $C$ depending only on $n$ and $\psi$, and uniform in $j$, $k$, $x_0$.
\end{proposition}

\subsection{Piece $T_1$: low--high paraproduct}

\begin{proposition}[Estimate for $T_1$]
\label{prop:T1}
Let $j \in [k, k+J]$.  Then
\[
  \norm{T_1 f}{L^2}
  \;\leq\; C_{J}\,\norm{m}_{L^\infty}
           \sum_{\ell : |\ell - j| \leq 3}
           \norm{S_{\ell-3}m}_{L^\infty}\,\norm{\Pi_\ell f}{L^2},
\]
with $C_J \leq C\cdot 2^{J/2}$ (bounded in $J$), and the sum has $O(1)$
terms.  Summing over $j \in [k, k+J]$ yields
\[
  \left(\sum_{j=k}^{k+J}\norm{T_1[\Pi_j] f}{L^2}^2\right)^{1/2}
  \;\leq\; C(J+1)^{1/2}\,\norm{m}_{W^{n/2+1,\infty}}\,\norm{f}{L^2}.
\]
\end{proposition}

\begin{proof}[Sketch]
By Fourier support considerations, $\Pi_j(\para_{(m-1)} f)$ is supported at
frequency $\sim 2^j$.  Since $\norm{m-1}_{L^\infty} \leq 1$ (as $m$ is a
smooth cutoff with values in $[0,1]$), the $L^2$ norm is controlled by
Bernstein's inequality.  The factor $(J+1)^{1/2}$ comes from summing $J+1$
terms.
\end{proof}

\subsection{Piece $T_2$: high--low paraproduct}

\begin{proposition}[Estimate for $T_2$]
\label{prop:T2}
For $j \in [k, k+J]$, the multiplier $m = \chi_{x_0,k}$ satisfies
$\Pi_\ell m = 0$ unless $|\xi| \lesssim 2^k$, so $T_2$ is non-trivial
only for $\ell \leq k + C_0$ (some absolute constant $C_0$).
Hence for $j \in [k, k+J]$ with $J$ fixed:
\[
  \norm{T_2 f}{L^2}
  \;\leq\; C\,\norm{\nabla^n \chi}_{L^1}
           \sum_{\ell=k}^{k+J+C_0}
           \norm{S_{\ell-3}f}_{L^\infty}\,\norm{\Pi_\ell m}_{L^2}.
\]
Applying Bernstein ($\norm{\Pi_\ell m}_{L^2} \leq C 2^{-\ell n/2}\norm{m}_{L^1}
\leq C 2^{-\ell n/2} \cdot 2^{-kn}$) and Sobolev embedding, we obtain
the $L^2$ bound uniform in $x_0$:
\[
  \norm{T_2 f}{L^2}
  \;\leq\; C_J\,\norm{\chi}_{H^{n/2+1}}\,\norm{f}{L^2}.
\]
The constant $C_J$ grows at most polynomially in $J$.
\end{proposition}

\subsection{Piece $T_3$: resonant/diagonal term}

\begin{proposition}[Estimate for $T_3$]
\label{prop:T3}
For $j \in [k, k+J]$,
\[
  T_3 f \;=\; m\,\Pi_j\!\left(\sum_{\ell}\Pi_\ell(m-1) \cdot \widetilde\Pi_\ell f\right).
\]
Since the spectrum of $(m-1)$ lies in $\{|\xi| \lesssim 2^k\}$ (as $m-1$
is also smooth at scale $2^{-k}$), the resonant
sum has contributions only from $\ell \in [k-C_0,\, k+J+C_0]$.  Thus
\[
  \norm{T_3 f}{L^2}
  \;\leq\; C\,\norm{m}_{L^\infty}
  \sum_{\ell=k-C_0}^{k+J+C_0}
  \norm{\Pi_\ell(m-1)}_{L^\infty}\,\norm{\Pi_j\widetilde\Pi_\ell f}{L^2}.
\]
Using $\norm{\Pi_\ell(m-1)}_{L^\infty} \leq C_\chi 2^{-k(N-n/2)}$ for any
$N \geq 0$ (since $m-1$ is smooth at scale $2^{-k}$), we get
\begin{equation}
  \label{eq:T3bound}
  \norm{T_3 f}{L^2}
  \;\leq\; C_{\chi,J}\, 2^{-k\varepsilon}\,\norm{f}{L^2}
\end{equation}
for any $\varepsilon > 0$, where $C_{\chi,J}$ depends only on $\chi$, $J$,
$n$, and $\varepsilon$, and is independent of $x_0$ and $k$.
\end{proposition}

\begin{remark}[Pointwise bound for $T_3$]
Under the additional assumption $f \in L^\infty$, the resonant piece
satisfies
\[
  \abs{T_3 f(x)}
  \;\leq\; C_{\chi,J}\,\norm{f}_{L^\infty}
           \sum_{\ell=k-C_0}^{k+J+C_0} 2^{\ell n}\,\norm{\Pi_\ell(m-1)}_{L^1}
  \;\leq\; C_{\chi,J}'\,(J+1)\,\norm{f}_{L^\infty},
\]
uniformly in $x_0$, $k$, and $x$.
\end{remark}

\subsection{Summary table}

\begin{center}
\renewcommand{\arraystretch}{1.4}
\begin{tabular}{|c|l|c|c|}
\hline
\textbf{Piece} & \textbf{Type} & \textbf{$L^2$ bound} & \textbf{$J$-dep.} \\
\hline
$T_1$ & Low--high paraproduct & $C(J+1)^{1/2}\norm{f}{L^2}$ & $O(\sqrt{J})$ \\
$T_2$ & High--low paraproduct & $C_J\norm{f}{L^2}$ & polynomial \\
$T_3$ & Resonant (diagonal)   & $C_{\chi,J}2^{-k\varepsilon}\norm{f}{L^2}$ & polynomial \\
$T_4$ & $(m-1)\Pi_j f$ cutoff & $C\norm{f}{L^2}$ & none \\
\hline
\end{tabular}
\end{center}

All constants are independent of $x_0 \in \R^n$, uniform in $k \geq 0$,
and uniform in $j \in [k, k+J]$.

% -----------------------------------------------------------------------
\section{Constant Tracking and Uniformity}
\label{sec:constants}
% -----------------------------------------------------------------------

\begin{theorem}[Main Decomposition Estimate]
\label{thm:main}
Let $n \geq 2$, $J \geq 1$, $k \geq 0$, $x_0 \in \R^n$, and
$j \in [k, k+J]$.  Then
\[
  \norm{\Pi_{j,x_0,k} f - \Pi_j f}{L^2}
  \;\leq\; \mathcal{A}(J,\chi,n)\,\norm{f}{L^2},
\]
where
\[
  \mathcal{A}(J,\chi,n)
  \;:=\; C(n)\bigl[(J+1)^{1/2}
         + C_{\chi}\,(J+1)^2\bigr]
  \;\leq\; C(n,\chi)\,(J+1)^2,
\]
and $C(n)$, $C_\chi$ depend only on dimension $n$ and the $C^{n+4}$ norm
of $\chi$, \emph{not} on $x_0$, $k$, or $j$.
\end{theorem}

\begin{proof}
Sum the estimates of Propositions~\ref{prop:T4}--\ref{prop:T3}.
The polynomial bound in $J$ absorbs all terms.
\end{proof}

\begin{corollary}[Uniformity in $(x_0, k)$]
The bound $\mathcal{A}(J,\chi,n)$ is uniform over all $x_0 \in \R^n$ and
$k \geq 0$.  In particular, the decomposition is suitable for insertion
into any summation over $k$ and spatial centres $x_0$ appearing in the
D9 envelope.
\end{corollary}

\begin{remark}[$J$-dependence is sharp for $T_3$]
\label{rem:Jsharp}
For the resonant piece $T_3$, the factor $(J+1)$ counts the number of
resonant frequency bands.  A simple example with $f = \sum_{\ell=k}^{k+J}
\Pi_\ell g$ and $m = $ a bump shows this is optimal up to constants.
\end{remark}

% -----------------------------------------------------------------------
\section{Off-Diagonal Ranges}
\label{sec:offdiag}
% -----------------------------------------------------------------------

\begin{remark}[Off-diagonal ranges]
\label{rem:offdiag}
The near-diagonal range $j \in [k, k+J]$ is the critical case.
\begin{itemize}
  \item For $j > k+J$ (sub-localisation scale): the projector $\Pi_j$ acts
    at frequencies below $2^k$, so $\Pi_{j,x_0,k} = \Pi_j$ on smooth
    data, and $\delta\Pi_{j,x_0,k} = 0$ up to exponentially small tails.
  \item For $j < k$ (super-localisation scale): $\Pi_j$ acts at frequencies
    $2^j \ll 2^k$, and the commutator $\comm(m,\Pi_j)$ has kernel decay
    $O(2^{j(N-n/2)} / 2^{kN})$ for any $N$, giving rapid decay in $(k-j)$.
\end{itemize}
Hence the only active range is indeed $j \in [k, k+J]$.
\end{remark}

% -----------------------------------------------------------------------
\section{Link to the D9 Envelope Framework}
\label{sec:d9link}
% -----------------------------------------------------------------------

The decomposition~\eqref{eq:main} and Theorem~\ref{thm:main} are designed
for direct insertion into the \emph{NS D9 envelope framework}, where one
constructs energy envelopes $\mathcal{E}_{k}(t)$ of the form
\[
  \mathcal{E}_{k}(t)
  \;:=\; \sup_{x_0}\sum_{j=k}^{k+J}
         \norm{\Pi_{j,x_0,k} u(t,\cdot)}{L^2(B(x_0,2^{-k}))}^2.
\]
The difference $\Pi_{j,x_0,k} - \Pi_j$ enters when commuting the
localisation with the nonlinear (Navier--Stokes) evolution.  The estimate
of Theorem~\ref{thm:main} shows that this commutator contributes at most
$\mathcal{A}(J,\chi,n)^2\,\norm{u}{L^2}^2$ to the energy envelope, with
constants that are uniform in $(x_0, k)$, enabling a closed estimate.

\paragraph{Insertion recipe.}
Given a D9-envelope inequality of the form
\[
  \partial_t \mathcal{E}_k
  \;\leq\; [\text{NS terms}]
           + R_k
\]
where $R_k$ involves $\Pi_{j,x_0,k} u - \Pi_j u$, replace $R_k$ by
\[
  R_k \;\leq\; \mathcal{A}(J,\chi,n)^2 \norm{u(t)}{L^2}^2
  \;+\; \text{(off-diagonal tails from Section~\ref{sec:offdiag})}.
\]
The off-diagonal tails are summable in $k$ with the factor
$\sum_{k \geq 0} 2^{-k\varepsilon} < \infty$.

% -----------------------------------------------------------------------
\section{Open Items and Next Steps}
\label{sec:open}
% -----------------------------------------------------------------------

\begin{itemize}
  \item \textbf{Sharp constants for $T_2$, $T_3$:}  The polynomial bound
    in $J$ for $T_2$ and $T_3$ may be improvable to $(J+1)^{1/2}$
    by bilinear estimates or $T(b)$-theorem techniques.
  \item \textbf{$L^p$ extensions:}  The estimates above are stated in
    $L^2$; extensions to $L^p$, $1 < p < \infty$, follow from
    Calderón--Zygmund theory but require tracking $p$-dependent constants.
  \item \textbf{Sobolev-norm version:}  For NS regularity theory one needs
    the analogue with $H^s$ norms; the same decomposition applies with
    $2^{sj}$ weights inserted.
  \item \textbf{N9.2 interface:}  The present decomposition feeds into
    the N9.2 \#8 observable definition once the admissibility conditions
    and windowing parameters are fixed.
\end{itemize}

% -----------------------------------------------------------------------
\appendix
\section{Bernstein and Littlewood--Paley Reference Inequalities}
\label{sec:bernstein}
% -----------------------------------------------------------------------

For reference, we collect the standard inequalities used above.

\begin{lemma}[Bernstein]
Let $f$ be supported at frequencies $|\xi| \leq R$.  Then for
$1 \leq p \leq q \leq \infty$:
\[
  \norm{D^\alpha f}{L^q}
  \;\leq\; C_\alpha\, R^{|\alpha| + n(1/p - 1/q)}\,\norm{f}{L^p}.
\]
\end{lemma}

\begin{lemma}[Littlewood--Paley square function]
For $1 < p < \infty$ there exists $C_p > 0$ such that
\[
  C_p^{-1}\,\norm{f}{L^p}
  \;\leq\; \norm{\Bigl(\sum_j |\Pi_j f|^2\Bigr)^{1/2}}{L^p}
  \;\leq\; C_p\,\norm{f}{L^p}.
\]
\end{lemma}

\begin{lemma}[Paraproduct $L^2$ bound]
For $a \in L^\infty$ and $f \in L^2$:
\[
  \norm{\para_a f}{L^2} \;\leq\; C\,\norm{a}_{L^\infty}\,\norm{f}{L^2},
  \qquad
  \norm{\res(a,f)}{L^2} \;\leq\; C\,\norm{a}_{L^\infty}\,\norm{f}_{L^2}.
\]
\end{lemma}

% -----------------------------------------------------------------------
\end{document}
